\documentclass[11pt,a4paper]{article}
\usepackage[utf8]{inputenc}
\usepackage[margin=0.4in]{geometry}
\usepackage{amsmath}
\usepackage{booktabs}
\usepackage{listings}
\usepackage{xcolor}
\usepackage{multicol}
\usepackage{graphicx}
\usepackage{array}

% Code highlighting configuration
\definecolor{mDarkTeal}{rgb}{0.0, 0.4, 0.4}
\definecolor{mGray}{rgb}{0.5,0.5,0.5}
\definecolor{mPurple}{rgb}{0.58,0,0.82}
\definecolor{backgroundHex}{rgb}{0.97,0.97,0.97}

\lstset{
    backgroundcolor=\color{backgroundHex},   
    commentstyle=\color{mDarkTeal}\itshape,
    keywordstyle=\color{magenta},
    numberstyle=\tiny\color{mGray},
    stringstyle=\color{mPurple},
    basicstyle=\ttfamily\scriptsize,
    breakatwhitespace=false,         
    breaklines=true,                 
    captionpos=b,                    
    keepspaces=true,                 
    numbers=none,                    
    showspaces=false,                
    showstringspaces=false,
    showtabs=false,                  
    tabsize=2,
    frame=single,
    aboveskip=6pt,
    %belowskip=2pt
}

\title{\vspace{-0.6in}\textbf{GATE Question Analysis \& Hardware Implementation: Full Adder}\vspace{-0.6in}}
%\author{\textbf{Technical Documentation Lab Reference}}
\date{}

\begin{document}

\maketitle
\thispagestyle{empty}

% --- FULL WIDTH SECTION: QUESTION ---
\noindent\makebox[\linewidth]{\rule{\textwidth}{0.4pt}}
\subsection*{The GATE Question}
\noindent Fig. 1 shows the $i^{th}$ full-adder block of a binary adder circuit. $C_i$ is the input carry and $C_{i+1}$ is the output carry of the circuit. If the inputs $A_i, B_i$ are available and stable throughout the carry propagation, find the outputs $S_i$ and $C_{i+1}$.

\begin{figure}[h!]
    \centering
    % If you have the image saved locally as an image file (e.g., gate_circuit.png), 
    % uncomment the line below and add your filename:
     \includegraphics[width=0.45\textwidth]{fa1.png}
    
    % This acts as a perfect visual placeholder matching your page layout requirements:
    %\framebox{\begin{minipage}{0.45\textwidth}
        \centering\vspace{0.15in}
        %\textbf{Fig. 3.13: Full Adder Logic Circuit Diagram} \\
        %\small Inputs: $A_i, B_i, C_i$ \ \ \textbar \ \ Outputs: $S_i, C_{i+1}$ \\
        \vspace{0.0000001in}
    %\end{minipage}}
\end{figure}
\noindent\makebox[\linewidth]{\rule{\textwidth}{0.4pt}}

\vspace{0.05in}

% --- TWO COLUMN SECTION: SOLUTION & LAB DATA ---
\begin{multicols}{2}

\section{Circuit Analytical Solution}
The given logical architecture is a classic standard Full Adder. By tracing the logic gate layout from the inputs, we derive the Boolean behavioral equations:
\begin{align*}
S_i &= A_i \oplus B_i \oplus C_i \\
C_{i+1} &= (A_i \cdot B_i) + (B_i \cdot C_i) + (A_i \cdot C_i)
\end{align*}

\section{Laboratory Components Used}
\begin{itemize}
    \setlength{\itemsep}{0pt}
    \item 1 $\times$ Arduino Uno Development Board
    \item 3 $\times$ Tactile Push buttons (Active-LOW style)
    \item 2 $\times$ LEDs (Red: $S_i$ Output, Red: $C_{i+1}$ Output)
    \item 2 $\times$ $220\Omega$ Current-Limiting Resistors
    \item 1 $\times$ Breadboard \& M-M Jumper Array
\end{itemize}

\section{Hardware Connection Mapping}
\begin{itemize}
    \setlength{\itemsep}{0pt}
    \item \textbf{Inputs:} $A_i \rightarrow$ Pin 2, $B_i \rightarrow$ Pin 3, $C_i \rightarrow$ Pin 4. All pins are bridged to \textbf{GND} via buttons using internal pull-up logic.
    \item \textbf{Outputs:} $S_i \rightarrow$ Pin 8 $\rightarrow 220\Omega \rightarrow$ Red LED. \\ $C_{i+1} \rightarrow$ Pin 9 $\rightarrow 220\Omega \rightarrow$ Green LED.
\end{itemize}

\vfill\columnbreak

\section{System Truth Table}
\centering
\small
\begin{tabular}{ccccc}
\toprule
$\mathbf{A_i}$ & $\mathbf{B_i}$ & $\mathbf{C_i}$ & \textbf{Sum ($\mathbf{S_i}$)} & \textbf{Carry ($\mathbf{C_{i+1}}$)} \\
\midrule
0 & 0 & 0 & 0 & 0 \\
1 & 0 & 0 & 1 & 0 \\
0 & 1 & 0 & 1 & 0 \\
1 & 1 & 0 & 0 & 1 \\
0 & 0 & 1 & 1 & 0 \\
1 & 0 & 1 & 0 & 1 \\
0 & 1 & 1 & 0 & 1 \\
1 & 1 & 1 & 1 & 1 \\
\bottomrule
\end{tabular}

\end{multicols}

\vspace{-0.8in}
\pagebreak

\setlength{\columnsep}{0.9cm}

\begin{multicols}{2}

\section{Arduino Embedded Cpp Sketch}
\begin{lstlisting}[language=C++]
const int pinA = 2, pinB = 3, pinCin = 4;
const int pinSum = 8, pinCout = 9;

void setup() {
  pinMode(pinA, INPUT_PULLUP);
  pinMode(pinB, INPUT_PULLUP);
  pinMode(pinCin, INPUT_PULLUP);
  pinMode(pinSum, OUTPUT);
  pinMode(pinCout, OUTPUT);
}

void loop() {
  // Read and invert for active-low buttons
  int A = !digitalRead(pinA);
  int B = !digitalRead(pinB);
  int Cin = !digitalRead(pinCin);

  // Digital execution
  int sum = A ^ B ^ Cin; 
  int cout = (A&&B) || (B&&Cin) || (A&&Cin);

  digitalWrite(pinSum, sum);
  digitalWrite(pinCout, cout);
  delay(50);
}
\end{lstlisting}

\vfill\columnbreak

\section{Assembly Language}
\begin{lstlisting}[language=Verilog]
; Target Register & I/O Address Definitions
.device ATmega328P
.equ DDRD=0x0A  \ \ .equ PORTD=0x0B  \ \ .equ PIND=0x09
.equ DDRB=0x04  \ \ .equ PORTB=0x05

setup:
    cbi DDRD,2 \ \ cbi DDRD,3 \ \ cbi DDRD,4
    sbi PORTD,2 \ sbi PORTD,3 \ sbi PORTD,4
    sbi DDRB,0 \ \ sbi DDRB,1

loop:
    in   r16, PIND
    com  r16          ; Handle pull-up state
    
    mov  r17, r16     ; Isolate A
    lsr  r17 \ lsr r17 \ andi r17, 0x01
    
    mov  r18, r16     ; Isolate B
    lsr  r18 \ lsr r18 \ lsr r18 \ andi r18, 0x01
    
    mov  r19, r16     ; Isolate Cin
    lsr  r19 \ lsr r19 \ lsr r19 \ lsr r19 \ andi r19, 0x01

    mov  r20, r17     ; XOR Logic for Sum
    eor  r20, r18 \ eor r20, r19

    mov  r21, r17 \ and r21, r18  ; A & B
    mov  r22, r18 \ and r22, r19  ; B & Cin
    mov  r23, r17 \ and r23, r19  ; A & Cin
    or   r21, r22 \ or  r21, r23  ; Carry logic

    lsl  r21          ; Position Cout to PB1
    or   r20, r21     ; Complete Output Pack
    out  PORTB, r20   ; Drive System LEDs
    rjmp loop
\end{lstlisting}

\end{multicols}

\end{document}
